\chapter{How texweb works}\label{ch:intro}
These notes are written in \LaTeX{}, as usual: the same source produces the PDF with \texttt{pdflatex} and the website with \texttt{texweb dev}. References, citations and numbering work in both; for example, chapter~\ref{ch:conv} is about convolution and cites \cite{Goodfellow-et-al-2016}.

\section{What you can write}
The dialect covers everyday commands: \textbf{bold}, \emph{italics}, \texttt{code}, footnotes\footnote{Like this one, numbered per page.} and formulas such as $f(x) = \max\{0, x\}$ or
\begin{equation}\label{eq:loss}
  \mathcal{L}(\vartheta) = \frac{1}{N}\sum_{i=1}^{N} \big(y^{(i)} - h(x^{(i)}; \vartheta)\big)^2, \qquad \vartheta \in \R^d.
\end{equation}
Equation~\eqref{eq:loss} is the MSE.

\begin{definition}[Loss function]\label{def:loss}
A \textbf{loss function} measures the error the model makes when predicting $y$ from $x$.
\end{definition}

Definition~\ref{def:loss} is numbered within the chapter, as in the PDF.

\begin{itemize}
  \item bulleted lists,
  \item numbered lists and tables,
  \item TikZ figures, compiled by \LaTeX{} into SVG.
\end{itemize}

\begin{table}[htbp]
\centering
\begin{tabular}{@{}lrr@{}}
\toprule
Model & Parameters & Layers \\
\midrule
MLP & 397\,510 & 2 \\
CNN & 156 & 1 \\
\bottomrule
\end{tabular}
\caption{Parameters of the first layer on MNIST.}
\label{tab:param}
\end{table}

Table~\ref{tab:param} shows how many parameters are saved.
